% Cross-Benchmark Integrated Analysis: ARM Semantic Modes + Quantitative Entropy Signals
% Supports extension of gpqa_loop_analysis.tex findings to HLE and AIME
% Requires: \usepackage{tcolorbox} and \findingbox command from gpqa_macro_analysis.tex

\subsection{Cross-Benchmark ARM and Context Rotting Analysis}
\label{appendix:macro-cross}

The loop-mechanics analysis in \S\ref{appendix:macro-gpqa-loop} decomposes GPQA-Diamond
trajectory inflation at the cognitive-mode and entropy levels.
We now extend the ARM (Agentic Reasoning Mode) taxonomy and the context-rotting
diagnostics to HLE, and cross-reference semantic mode signatures with per-token
entropy signals across both benchmarks.
Four findings emerge that are invisible within a single-benchmark analysis.

% ===========================================================================
\paragraph{The timeout budget is filled by near-deterministic repetition, not
tool execution.}
% ===========================================================================

Table~\ref{tab:cross-timeout} decomposes OpenClaw trajectories on GPQA-Diamond
into three regimes: correct (mean 2.3K tokens), wrong-but-not-capped (10.0K),
and context-exhausted ($\ge$30K tokens before answer emission).
The capped regime reveals exactly what fills the additional 44K tokens.

\begin{table}[t]
\centering
\small
\caption{%
  Token budget composition for GPQA-Diamond OpenClaw trajectories, stratified
  by outcome and cap status.
  Capped = trajectories reaching $\ge$30K tokens before emitting a valid boxed answer.
  Arm entropy = cognitive mode diversity.
  Malformed\,\% = extracted answer does not match any valid option letter.
  No-boxed\,\% = no \texttt{\textbackslash boxed\{\}} pattern found.
}
\label{tab:cross-timeout}
\setlength{\tabcolsep}{2.5pt}
\begin{tabular}{lrrr}
\toprule
& Capped ($n{=}25$) & Wrong, not capped ($n{=}29$) & Correct ($n{=}131$) \\
\midrule
Mean tokens            & 53{,}887 & 10{,}035 & 2{,}319 \\
Mean entropy           & 0.030 & 0.218 & 0.248 \\
Head entropy           & 0.237 & 0.258 & 0.249 \\
Tail entropy           & \textbf{0.001} & 0.145 & 0.209 \\
$\Delta$H (tail $-$ head) & \textbf{$-$0.236} & $-$0.114 & $-$0.040 \\
Low-entropy token share & 0.955 & 0.706 & 0.621 \\
Repeated 4-gram rate   & 0.962 & 0.436 & 0.135 \\
Top-1 token probability & 0.990 & 0.924 & 0.909 \\
\midrule
Tool-intent markers    & 89.1 & 23.1 & 2.1 \\
Actual tool calls      & \textbf{0.0} & 0.2 & 0.4 \\
\midrule
Looping markers        & 188.2 & 23.9 & 2.3 \\
Self-correction markers & 283.2 & 43.0 & 5.9 \\
Uncertainty markers    & 90.8 & 4.8 & 1.2 \\
\midrule
ARM: PD rate           & 0.660 & 0.601 & 0.538 \\
ARM: reason rate       & 0.230 & 0.295 & 0.372 \\
ARM: IC rate           & \textbf{0.000} & 0.026 & 0.049 \\
ARM: UN rate           & 0.090 & 0.021 & 0.016 \\
ARM: RR rate           & \textbf{0.000} & 0.023 & 0.017 \\
Arm entropy            & 0.319 & 0.436 & 0.523 \\
Mean segments          & 2.1 & 4.3 & 4.7 \\
\midrule
Malformed\,\%          & 100 & 41 & 1 \\
No-boxed\,\%           & 100 & 34 & 1 \\
\bottomrule
\end{tabular}
\end{table}

\findingbox{Context-exhausted trajectories are not ``long reasoning'';
they are near-deterministic cognitive collapse.
The model enters a single PD segment, never exits, and its output entropy
approaches the minimum physically attainable value.}

The capped regime is not an extreme point on a continuum of increasing
deliberation; it is a phase transition.
Entropy enters at a normal 0.24 (head) and exits at 0.001 (tail), a drop of
0.24 nats compared to 0.04 for correct trajectories.
At 0.001 nats, the model's output is effectively a fixed string: 96\% of
all 4-grams are repeats, and the model assigns probability 0.99 to its
selected token.
The 89 tool-intent markers per trajectory express a desire to compute that
the model never fulfills: actual tool calls are zero.
The ARM taxonomy confirms the cognitive signature: the 2.1-segment trajectory
is a single PD monolith with IC rate zero and RR rate zero.
The model cannot commit and cannot recover; it can only re-derive.

Critically, non-capped wrong trajectories are not miniature versions of this.
They average 4.3 segments with IC at 0.03 and RR at 0.02: the model makes
commitment and recovery attempts, they just fail.
The capped regime is categorically different: the exit mechanisms are absent,
not merely ineffective.

% ===========================================================================
\paragraph{The PD-sink failure mode is benchmark-specific, not a universal
property of the model.}
% ===========================================================================

Table~\ref{tab:cross-arm-benchmark} extends the ARM taxonomy to HLE
(DirectLLM and OpenCode), revealing a fundamental shift in cognitive strategy
between the two benchmarks.

\begin{table}[t]
\centering
\small
\caption{%
  ARM cognitive-mode distribution by benchmark and outcome.
  GPQA-Diamond values from the ARM taxonomy described in
  \S\ref{appendix:macro-gpqa-loop}.
  HLE values from the same keyword classifiers applied to HLE trajectory text.
  $\Delta$PD = PD(wrong) $-$ PD(correct).
}
\label{tab:cross-arm-benchmark}
\setlength{\tabcolsep}{3pt}
\begin{tabular}{llrrrrrrr}
\toprule
Bench. & Harness & Out. & PD & reason & IC & UN & RR & $\Delta$PD \\
\midrule
\multirow{6}{*}{GPQA}
& \multirow{2}{*}{OpenClaw}
  & C & 0.538 & 0.372 & 0.049 & 0.016 & 0.017 & \multirow{2}{*}{+0.070} \\
& & W & 0.608 & 0.291 & 0.026 & 0.043 & 0.010 & \\
\cmidrule{2-8}
& \multirow{2}{*}{OpenCode}
  & C & 0.353 & 0.587 & 0.026 & 0.009 & 0.019 & \multirow{2}{*}{+0.145} \\
& & W & 0.498 & 0.449 & 0.000 & 0.042 & 0.011 & \\
\cmidrule{2-8}
& \multirow{2}{*}{ZeroClaw}
  & C & 0.280 & 0.664 & 0.023 & 0.011 & 0.021 & \multirow{2}{*}{+0.126} \\
& & W & 0.406 & 0.504 & 0.011 & 0.032 & 0.044 & \\
\midrule
\multirow{4}{*}{HLE}
& \multirow{2}{*}{DirectLLM}
  & C & 0.125 & 0.417 & 0.082 & 0.099 & 0.272 & \multirow{2}{*}{\textbf{$-$0.019}} \\
& & W & 0.106 & 0.423 & 0.074 & 0.135 & 0.259 & \\
\cmidrule{2-8}
& \multirow{2}{*}{OpenCode}
  & C & 0.162 & 0.142 & 0.132 & 0.201 & 0.364 & \multirow{2}{*}{\textbf{$-$0.017}} \\
& & W & 0.144 & 0.153 & 0.152 & 0.238 & 0.310 & \\
\bottomrule
\end{tabular}
\end{table}

\findingbox{The PD-sink over-reasoning signature from GPQA-Diamond does not
appear on HLE.
ARM failure modes are determined by the interaction between benchmark domain
structure and harness architecture, not by a fixed model-level cognitive bias.}

On GPQA-Diamond, principle-based deduction is the dominant mode across all
harnesses (28\% to 61\% of segments), and PD over-expression is the
strongest single-mode predictor of failure: the wrong$-$correct PD gap is
+7.0 to +14.5 percentage points across harnesses.
On HLE, PD is a minority mode (11\% to 16\%), and the wrong$-$correct PD gap
is \emph{negative} for both DirectLLM ($-$1.9~pp) and OpenCode ($-$1.7~pp).
Wrong HLE trajectories allocate \emph{less} of their cognitive budget to
deduction, not more.

Three structural differences drive the divergence.
First, GPQA-Diamond questions are hard-science problems where PD is the
natural reasoning strategy; HLE questions span philosophy, linguistics,
theoretical mathematics, and experimental science, where no single
derivational framework applies.
Second, the dominant non-reason mode on HLE is recovery (RR at 0.26 to
0.36), compared to $<$0.05 on GPQA: the model is constantly reconsidering
its approach under genuine domain uncertainty, not spiraling in a single
mode.
Third, arm entropy on HLE is substantially higher (1.14 to 1.23 for
DirectLLM) than on GPQA (0.40 to 0.81), confirming that mode diversity
reflects task diversity rather than failure.

The implication is that a harness designed to prevent the PD spiral on
GPQA (e.g., ZeroClaw's fixed template) may have no effect on HLE, where
the failure mode is not over-deduction but premature commitment under
knowledge uncertainty.

% ===========================================================================
\paragraph{Context rotting is harness-conditional and benchmark-conditional;
a catastrophic threshold exists below which answer extraction becomes
impossible.}
% ===========================================================================

Table~\ref{tab:cross-rotting} extends the context-rotting analysis from
\S\ref{appendix:macro-gpqa-loop} (GPQA only) to HLE, covering six
harness-benchmark combinations.

\begin{table}[t]
\centering
\small
\caption{%
  Context rotting signals across GPQA-Diamond and HLE.
  Head = mean entropy over first quartile of tokens.
  Tail = mean over last quartile.
  $\Delta$H = tail $-$ head.
  Low-ent\,\% = fraction of tokens in the bottom quartile of the
  per-trajectory entropy distribution.
  Top-1 = mean probability of the model's selected token.
}
\label{tab:cross-rotting}
\setlength{\tabcolsep}{3pt}
\begin{tabular}{llrrrrrrr}
\toprule
Bench. & Harness & Out. & Tokens & Head & Tail & $\Delta$H & Low-ent\,\% & Top-1 \\
\midrule
\multirow{6}{*}{GPQA}
& \multirow{2}{*}{OpenClaw}
  & C &  2{,}319 & 0.249 & 0.209 & $-$0.040 & 0.621 & 0.909 \\
& & W & 24{,}484 & 0.246 & 0.106 & \textbf{$-$0.140} & \textbf{0.808} & \textbf{0.952} \\
\cmidrule{2-9}
& \multirow{2}{*}{OpenCode}
  & C & 1{,}067 & 0.266 & 0.215 & $-$0.050 & 0.573 & 0.892 \\
& & W & 3{,}302 & 0.246 & 0.178 & $-$0.068 & 0.623 & 0.901 \\
\cmidrule{2-9}
& \multirow{2}{*}{ZeroClaw}
  & C & 2{,}102 & 0.236 & 0.183 & $-$0.053 & 0.635 & 0.911 \\
& & W & 2{,}338 & 0.244 & 0.254 & \textbf{+0.010} & \textbf{0.529} & \textbf{0.873} \\
\midrule
\multirow{6}{*}{HLE}
& \multirow{2}{*}{DirectLLM}
  & C & 12{,}175 & 0.402 & 0.271 & $-$0.131 & 0.249 & 0.879 \\
& & W & 14{,}710 & 0.407 & 0.260 & $-$0.147 & 0.244 & 0.882 \\
\cmidrule{2-9}
& \multirow{2}{*}{OpenCode}
  & C &  1{,}466 & 0.372 & 0.328 & $-$0.044 & 0.250 & 0.854 \\
& & W &  7{,}373 & 0.324 & 0.262 & $-$0.061 & 0.250 & 0.885 \\
\cmidrule{2-9}
& \multirow{2}{*}{ZeroClaw}
  & C &  2{,}340 & 0.334 & 0.272 & $-$0.062 & 0.250 & 0.880 \\
& & W &  2{,}287 & 0.338 & 0.271 & $-$0.067 & 0.250 & 0.879 \\
\bottomrule
\end{tabular}
\end{table}

\findingbox{Entropy decline per se is not a failure signature; only
catastrophic entropy collapse to near-zero values distinguishes
context-exhaustion failures from normal long-form deliberation.}

Table~\ref{tab:cross-rotting} reveals three distinct rotting regimes.

\textbf{Regime 1: No rotting.}
ZeroClaw maintains flat or rising entropy on both benchmarks.
On GPQA, its wrong trajectories show $\Delta$H = +0.01 (entropy \emph{rises}
from head to tail) and the lowest low-entropy token share (0.53) of any
wrong-trajectory cell.
On HLE, ZeroClaw shows uniform $\Delta$H $\approx$ $-$0.06 across outcomes.
The fixed template prevents decay: the trajectory ends before the rotting
regime begins, and the model remains in a deliberative state throughout.

\textbf{Regime 2: Outcome-invariant decline.}
HLE DirectLLM trajectories lose approximately 0.14 nats from head to tail
regardless of correctness.
Both correct and wrong trajectories start at approximately 0.40 and end at
approximately 0.26.
This is a natural consequence of long-form reasoning: as the model narrows
toward a conclusion, its distribution sharpens.
The decline is not diagnostic of failure.

\textbf{Regime 3: Catastrophic collapse.}
GPQA OpenClaw wrong trajectories drop from 0.25 to 0.11 ($\Delta$H = $-$0.14),
and the capped subset (Table~\ref{tab:cross-timeout}) plunges to 0.001.
The low-entropy token share rises to 0.81 (vs.\ 0.62 for correct), and top-1
probability reaches 0.95.
At these levels, the model's output is functionally deterministic: it cycles
through the same high-confidence restatements without advancing toward a
decision.

The critical threshold appears to be at tail entropy approximately 0.05.
Above this threshold (HLE DirectLLM at 0.26, OpenCode wrong at 0.18),
answer extraction is reliable.
Below it (OpenClaw capped at 0.001), 100\% of trajectories are malformed
and 100\% lack a boxed answer.

% ===========================================================================
\paragraph{Recovery (RR) is not a portable quality signal: its diagnostic
meaning depends jointly on benchmark domain structure and arm entropy.}
% ===========================================================================

Table~\ref{tab:cross-rr} reports RR rates alongside arm entropy and PD rate
across outcomes and benchmarks.

\begin{table}[t]
\centering
\small
\caption{%
  Recovery (RR) mode rates in context.
  Arm entropy = cognitive mode diversity (higher = more modes used within a
  trajectory).
  PD = principle-based deduction rate.
  GPQA outcomes use the four-way paired decomposition from
  \S\ref{appendix:macro-gpqa}.
}
\label{tab:cross-rr}
\setlength{\tabcolsep}{3pt}
\begin{tabular}{llrrrrr}
\toprule
Bench. & Harness & Outcome & RR rate & Arm ent. & PD rate & Tokens \\
\midrule
\multirow{12}{*}{GPQA}
& \multirow{4}{*}{OpenClaw}
  & both correct & 0.018 & 0.476 & 0.519 &  1{,}932 \\
& & both wrong   & 0.016 & 0.324 & 0.568 & 22{,}148 \\
& & regression   & 0.005 & 0.261 & 0.643 & 30{,}451 \\
& & rescue       & 0.000 & 0.161 & 0.831 &  3{,}598 \\
\cmidrule{2-7}
& \multirow{4}{*}{OpenCode}
  & both correct & 0.020 & 0.563 & 0.335 &    967 \\
& & both wrong   & 0.002 & 0.486 & 0.471 &  3{,}251 \\
& & regression   & 0.021 & 0.407 & 0.530 &  3{,}672 \\
& & rescue       & 0.000 & 0.386 & 0.600 &    701 \\
\cmidrule{2-7}
& \multirow{4}{*}{ZeroClaw}
  & both correct & 0.020 & 0.650 & 0.271 &  1{,}982 \\
& & both wrong   & 0.044 & 0.762 & 0.392 &  2{,}482 \\
& & regression   & 0.044 & 0.825 & 0.427 &  1{,}913 \\
& & rescue       & 0.032 & 0.679 & 0.380 &  2{,}333 \\
\midrule
\multirow{4}{*}{HLE}
& \multirow{2}{*}{DirectLLM}
  & correct & 0.272 & 1.230 & 0.125 & 12{,}175 \\
& & wrong   & 0.259 & 1.136 & 0.106 & 14{,}710 \\
\cmidrule{2-7}
& \multirow{2}{*}{OpenCode}
  & correct & 0.364 & 0.166 & 0.162 &  1{,}466 \\
& & wrong   & 0.310 & 0.128 & 0.144 &  7{,}373 \\
\bottomrule
\end{tabular}
\end{table}

\findingbox{The same RR rate carries opposite diagnostic meaning on different
benchmarks: it marks productive self-correction on HLE and mode oscillation on
GPQA ZeroClaw.}

On HLE, RR is the dominant non-reason cognitive mode at rates of 0.26 to
0.36, roughly 10$\times$ the GPQA baseline.
Arm entropy is high (1.14 to 1.23) and PD rates are low (0.11 to 0.16).
This profile describes a model that is genuinely navigating uncertainty: it
frequently reconsiders its approach, draws on diverse cognitive strategies,
and does not default to deduction under ambiguity.

On GPQA, RR rates are uniformly below 0.05, but the RR signal splits by
harness architecture.
Under OpenClaw and OpenCode, RR appears in 9\%--12\% of correct vs.\
3\%--4\% of wrong trajectories: recovery is a success signal.
Under ZeroClaw, the direction inverts: 49\% of wrong trajectories contain RR
vs.\ 24\% of correct ones.
But ZeroClaw wrong trajectories also carry the highest arm entropy of any
GPQA cell (0.81 for regression vs.\ 0.64 for correct): the RR signal is
embedded in mode oscillation rather than focused correction.

The quadruple (RR rate, arm entropy, PD rate, benchmark domain) jointly
determines whether recovery is productive.
High RR with high arm entropy and low PD = productive exploration (HLE).
High RR with high arm entropy and high PD = mode oscillation (ZeroClaw wrong).
Low RR with low arm entropy and high PD = cognitive collapse (OpenClaw capped).

% ===========================================================================
\paragraph{PD self-transition probability is the single strongest
quantitative predictor of context exhaustion.}
% ===========================================================================

Table~\ref{tab:cross-pd-transition} reports ARM transition probabilities
across all three GPQA harnesses.

\begin{table}[t]
\centering
\small
\caption{%
  PD-related ARM transition probabilities by harness and outcome on
  GPQA-Diamond.
  PD$\to$PD = probability that a PD segment is followed by another PD segment.
  reason$\to$PD = probability that undifferentiated reasoning transitions into PD.
  reason$\to$IC = probability that reasoning transitions to intuitive commitment
  (the natural exit from deliberation).
}
\label{tab:cross-pd-transition}
\setlength{\tabcolsep}{4pt}
\begin{tabular}{llrrrr}
\toprule
Harness & Outcome & PD$\to$PD & reason$\to$PD & reason$\to$IC & PD$\to$reason \\
\midrule
\multirow{2}{*}{OpenClaw}
  & correct & 0.729 & 0.513 & \textbf{0.356} & 0.248 \\
  & wrong   & \textbf{0.863} & \textbf{0.765} & \textbf{0.091} & 0.345 \\
\cmidrule{1-5}
\multirow{2}{*}{OpenCode}
  & correct & 0.251 & 0.301 & \textbf{0.198} & 0.222 \\
  & wrong   & \textbf{0.536} & \textbf{0.525} & \textbf{0.016} & 0.270 \\
\cmidrule{1-5}
\multirow{2}{*}{ZeroClaw}
  & correct & 0.215 & 0.165 & 0.079 & 0.163 \\
  & wrong   & 0.241 & 0.182 & 0.083 & 0.182 \\
\bottomrule
\end{tabular}
\end{table}

Three structural signatures distinguish the spiraling harness from the
non-spiraling ones.

First, the PD$\to$PD self-transition probability in OpenClaw wrong
trajectories is 0.86, meaning that once the model enters principle-based
deduction under uncertainty, the probability of remaining in deduction on the
next cognitive segment is 86\%.
This is a near-absorbing state.

Second, the commitment exit path collapses.
OpenClaw's reason$\to$IC transition drops from 0.36 (correct) to 0.09
(wrong): the model is 4$\times$ less likely to commit after reasoning.
OpenCode shows an even steeper collapse from 0.20 to 0.02 (10$\times$).
ZeroClaw's reason$\to$IC is uniformly low (0.08) regardless of outcome:
commitment is always forced by the template, not by cognitive conviction.

Third, the PD entry probability from reason (reason$\to$PD) rises sharply
on wrong trajectories for OpenClaw (0.51 to 0.77) and OpenCode (0.30 to
0.53), but not for ZeroClaw (0.17 to 0.18).
Under open-loop harnesses, uncertainty triggers re-entry into deduction;
under ZeroClaw's fixed template, the next step is predetermined and the
model cannot re-enter PD even when uncertain.

These three transition probabilities encode the mechanism of the PD sink:
a near-absorbing PD state, a collapsing exit path to commitment, and an
architecture-dependent re-entry probability from generic reasoning into
deduction.
A harness that breaks any one of the three prevents the spiral.
ZeroClaw breaks the re-entry probability by fixing step sequences.
OpenCode partially breaks the PD self-transition by providing tool-grounded
exit ramps.
Neither breaks all three, and both have wrong-trajectory PD rates elevated
by 12 to 15~pp over correct trajectories, but neither produces the
catastrophic entropy collapse of OpenClaw.

% ===========================================================================
\paragraph{Tool-type composition and quality: what tools each harness invokes
and whether they work.}
% ===========================================================================

The canonical action analysis tells us \emph{whether} a harness calls tools.
Here we decompose \emph{what kind} of tools and \emph{whether the calls are
substantive}.
OpenClaw provides a general-purpose tool shed (shell, read, search, write,
memory\_read); OpenCode provides a single code-execution environment;
ZeroClaw provides none.
Table~\ref{tab:cross-tool-types} reports tool-type distribution and quality
metrics by harness and outcome on GPQA-Diamond.

\begin{table}[t]
\centering
\small
\caption{%
  Tool-type composition and call quality by harness and outcome on
  GPQA-Diamond.
  Tool types are from event-log subtype labels.
  Any call = mean tool invocations per trajectory (including errors and
  empty outputs).
  Substantive = mean tool invocations producing executable code or
  non-empty output.
  Tool intent = keyword markers expressing desire to use a tool without
  necessarily executing one.
}
\label{tab:cross-tool-types}
\setlength{\tabcolsep}{3pt}
\begin{tabular}{llrrrrrr}
\toprule
Harness & Out. & Shell & Read & Search & Other & Any call & Subst. & Intent \\
\midrule
\multirow{2}{*}{OpenClaw}
  & C & 67\% & 12\% & 9\% & 12\% & 1.34 & 0.89 & 2.1 \\
  & W & 50\% & 17\% & 17\% & 17\% & 0.69 & 0.28 & \textbf{48.0} \\
\midrule
\multirow{2}{*}{OpenCode}
  & C & 0\% & 0\% & 0\% & 100\% & 0.95 & 0.95 & 7.8 \\
  & W & 0\% & 0\% & 0\% & 100\% & 0.32 & 0.32 & 11.4 \\
\midrule
\multirow{2}{*}{ZeroClaw}
  & C & \multicolumn{3}{c}{no tools} & & 0.00 & 0.00 & 6.8 \\
  & W & \multicolumn{3}{c}{no tools} & & 0.00 & 0.00 & 7.2 \\
\bottomrule
\end{tabular}
\end{table}

\findingbox{Tool-intent inflation without tool execution is the signature of
the OpenClaw spiral: intent markers explode 23-fold on wrong trajectories
while actual tool calls drop to near zero.
OpenCode's code-only tool surface produces the inverse pattern: every call
is substantive, but the call rate collapses when the model is uncertain.}

Three structural facts emerge from Table~\ref{tab:cross-tool-types}.

First, OpenClaw's tool surface is diverse but under-utilized.
Correct trajectories average 1.34 tool calls with 0.89 substantive (67\%
execution rate); wrong trajectories drop to 0.69 calls with 0.28 substantive
(41\% execution rate).
Simultaneously, tool-intent markers explode from 2.1 (correct) to 48.0
(wrong).
The model repeatedly expresses the desire to compute something
(\texttt{Let me compute the density}, \texttt{I need to check this
calculation}) but neither executes nor integrates the result.
The tool shed provides multiple exit paths (shell for computation, search
for fact-checking, read for reference), but under the PD spiral the model
expresses intent without acting on any of them.

Second, OpenCode's single-tool surface produces an all-or-nothing pattern.
Every tool call is substantive (code execution always produces output), but
the call rate drops 3.0$\times$ from correct to wrong (0.95 to 0.32).
The model's tool activation is outcome-predictive not because tools cause
correctness, but because the model only persists in calling tools when it is
on a productive reasoning path.
On wrong trajectories, the model abandons tool use early and defaults to
pure reasoning, producing the elevated PD rate documented in
Table~\ref{tab:cross-arm-benchmark}.

Third, ZeroClaw has no tool surface at all, yet its tool-intent markers
(6.8 to 7.2) are comparable to OpenCode's correct trajectories (7.8).
The model expresses approximately the same baseline level of computational
intent regardless of harness; the difference is that ZeroClaw's template
suppresses the expression at the action level, while OpenClaw's open loop
amplifies it into a repeating marker without execution.

% ===========================================================================
\paragraph{Action-level trajectory motifs: what operational patterns
distinguish correct from wrong trajectories.}
% ===========================================================================

Canonical actions $\{\texttt{answer}, \texttt{reason}, \texttt{tool\_use},
\texttt{verify}, \texttt{plan}, \texttt{recover}\}$ classify the operational
structure of each trajectory step.
Table~\ref{tab:cross-action-motifs} reports the top 2-gram and 3-gram action
motifs by harness and outcome on GPQA-Diamond.
A motif's rate is the mean number of occurrences per trajectory.

\begin{table}[t]
\centering
\small
\caption{%
  Top action motifs (2-gram and 3-gram transitions) by harness and outcome
  on GPQA-Diamond.
  Rate = mean occurrences per trajectory in the outcome bucket.
  Motifs are ordered by rate within each (harness, outcome) cell.
}
\label{tab:cross-action-motifs}
\setlength{\tabcolsep}{3pt}
\begin{tabular}{llllr}
\toprule
Harness & Outcome & \multicolumn{2}{c}{Top 3 motifs} & Rate \\
\midrule
\multirow{3}{*}{OpenClaw}
  & \multirow{3}{*}{C}
    & answer$\to$verify            & & 0.53 \\
  & & answer$\to$answer            & & 0.22 \\
  & & answer$\to$verify$\to$recover &  & 0.15 \\
\cmidrule{2-5}
  & \multirow{3}{*}{W}
    & answer$\to$verify            & & 0.45 \\
  & & answer$\to$plan              & \textbf{(re-planning)} & 0.27 \\
  & & recover$\to$recover          & \textbf{(oscillation)} & 0.19 \\
\midrule
\multirow{3}{*}{OpenCode}
  & \multirow{3}{*}{C}
    & answer$\to$reason            & & 0.86 \\
  & & reason$\to$verify            & & 0.49 \\
  & & tool\_use$\to$reason         & \textbf{(tool-grounded)} & 0.42 \\
\cmidrule{2-5}
  & \multirow{3}{*}{W}
    & answer$\to$reason            & & 0.92 \\
  & & reason$\to$verify            & & 0.46 \\
  & & reason$\to$plan              & \textbf{(stalled planning)} & 0.23 \\
\midrule
\multirow{3}{*}{ZeroClaw}
  & \multirow{3}{*}{C}
    & reason$\to$reason            & & 11.43 \\
  & & reason$\to$reason$\to$reason & & 9.77 \\
  & & plan$\to$reason              & & 1.23 \\
\cmidrule{2-5}
  & \multirow{3}{*}{W}
    & reason$\to$reason            & & 10.88 \\
  & & reason$\to$reason$\to$reason & & 9.03 \\
  & & reason$\to$plan              & \textbf{(diverted planning)} & 1.44 \\
\bottomrule
\end{tabular}
\end{table}

\findingbox{The dominant action motif on OpenClaw wrong trajectories shifts
from productive verification (answer$\to$verify) to re-planning without
execution (answer$\to$plan) and recovery oscillation
(recover$\to$recover).
OpenCode wrong trajectories lose the tool-grounded reasoning cycle
(tool\_use$\to$reason) that characterizes its correct trajectories.}

Table~\ref{tab:cross-action-motifs} reveals three harness-specific
operational signatures.

OpenClaw correct trajectories follow a verify-then-recover cycle:
answer$\to$verify at 0.53, with answer$\to$verify$\to$recover at 0.15.
The model answers, checks its work, and (when needed) corrects.
OpenClaw wrong trajectories retain answer$\to$verify at 0.45 but the
second-rank motif shifts to answer$\to$plan at 0.27: the model stops
verifying and starts re-planning without executing.
The third-rank motif becomes recover$\to$recover at 0.19: recovery
attempts chain into each other without resolution, the action-level
analogue of the ARM-level PD$\to$PD sink.

OpenCode correct trajectories are characterized by the tool\_use$\to$reason
motif at 0.42: the model executes code and reasons about the output.
On wrong trajectories, this motif drops below the top 3, replaced by
reason$\to$plan at 0.23.
The model still reasons and verifies (reason$\to$verify at 0.46, comparable
to 0.49 on correct), but it has lost the tool-grounded anchor and
substitutes planning for execution.

ZeroClaw trajectories are dominated by reason$\to$reason chains at rates
approximately 10 per trajectory regardless of outcome.
This is the template signature: the fixed plan-reason-answer structure
produces long contiguous reasoning blocks.
The only outcome-conditioned signal is a slight shift from plan$\to$reason
(correct, 1.23) to reason$\to$plan (wrong, 1.44): wrong trajectories
interrupt reasoning with planning, while correct ones feed planning into
reasoning.

% ===========================================================================
\paragraph{Action rates across benchmarks: tool use as an
outcome-discriminative signal is benchmark-conditional.}
% ===========================================================================

Table~\ref{tab:cross-action-rates} reports canonical action rates for all
harnesses on both GPQA-Diamond and HLE.
An action rate is the fraction of trajectories in the outcome bucket that
contain at least one event of that action type.

\begin{table}[t]
\centering
\small
\caption{%
  Canonical action rates by benchmark, harness, and outcome.
  Each value = fraction of trajectories in the (harness, outcome) cell
  containing $\ge$1 event of that action type.
  GPQA values from event-log extraction; HLE values from keyword heuristics
  on trajectory text (DirectLLM) and normalized trace step types (OpenCode,
  ZeroClaw).
}
\label{tab:cross-action-rates}
\setlength{\tabcolsep}{2.5pt}
\begin{tabular}{lllrrrrrr}
\toprule
Bench. & Harness & Out. & Ans. & Reas. & Tool & Verify & Plan & Recov. \\
\midrule
\multirow{14}{*}{GPQA}
& \multirow{2}{*}{DirectLLM}
  & C & 1.00 & 1.00 & 0.00 & 0.00 & 0.00 & 0.00 \\
& & W & 1.00 & 1.00 & 0.00 & 0.00 & 0.00 & 0.00 \\
\cmidrule{2-8}
& \multirow{2}{*}{OpenClaw}
  & C & 0.99 & 0.12 & 0.22 & 0.57 & 0.16 & \textbf{0.24} \\
& & W & 0.85 & 0.19 & \textbf{0.07} & 0.58 & \textbf{0.31} & 0.10 \\
\cmidrule{2-8}
& \multirow{2}{*}{OpenCode}
  & C & 1.00 & 0.93 & \textbf{0.55} & 0.59 & 0.30 & 0.00 \\
& & W & 1.00 & 0.92 & \textbf{0.15} & 0.49 & 0.32 & 0.00 \\
\cmidrule{2-8}
& \multirow{2}{*}{ZeroClaw}
  & C & 1.00 & 1.00 & 0.00 & 0.60 & 0.82 & 0.00 \\
& & W & 1.00 & 1.00 & 0.00 & 0.59 & 0.81 & 0.00 \\
\midrule
\multirow{12}{*}{HLE}
& \multirow{2}{*}{DirectLLM}
  & C & 1.00 & 1.00 & 0.00 & 0.99 & 0.58 & 1.00 \\
& & W & 0.86 & 1.00 & 0.01 & 0.97 & 0.47 & 0.99 \\
\cmidrule{2-8}
& \multirow{2}{*}{OpenCode}
  & C & 1.00 & 1.00 & \textbf{0.23} & 1.00 & 0.56 & 0.70 \\
& & W & 1.00 & 1.00 & \textbf{0.25} & 1.00 & 0.62 & 0.69 \\
\cmidrule{2-8}
& \multirow{2}{*}{ZeroClaw}
  & C & 1.00 & 1.00 & 0.00 & 1.00 & 0.73 & 0.73 \\
& & W & 1.00 & 1.00 & 0.00 & 1.00 & 0.65 & 0.71 \\
\bottomrule
\end{tabular}
\end{table}

\findingbox{The tool-use rate gap between correct and wrong trajectories is
large on GPQA-Diamond (3.7$\times$ for OpenCode) and near zero on HLE
(0.9$\times$ for OpenCode).
Tool activation is an outcome-discriminative signal only when the benchmark
contains executable subproblems.}

The GPQA rows of Table~\ref{tab:cross-action-rates} show a sharp
outcome-conditioned tool-use gap.
OpenCode invokes tools in 55\% of correct vs.\ 15\% of wrong trajectories:
a 3.7$\times$ difference.
OpenClaw shows a similar direction at 22\% vs.\ 7\% (3.1$\times$).
On HLE, OpenCode's tool-use rate is 23\% correct vs.\ 25\% wrong: the
outcome-conditioned gap collapses to 0.9$\times$.
The model invokes tools at the same rate regardless of whether it
ultimately answers correctly because HLE tool calls are exploratory
(fact-checking, information retrieval) rather than computational.

Two additional patterns are visible only in the cross-benchmark view.
First, recovery (Recov.) rates are uniformly higher on HLE than GPQA for
every harness.
OpenCode's recovery rate rises from 0.00 (GPQA) to 0.70 (HLE); ZeroClaw's
from 0.00 to 0.71--0.73.
This reflects the ARM-level finding that RR is the dominant non-reason mode
on HLE: the action-level data confirms that the model reconsiders and
corrects at the operational level.
Second, planning (Plan) rates are compressed toward uniformity on HLE:
the GPQA OpenClaw plan explosion (0.16 correct $\to$ 0.31 wrong) does not
appear on HLE, where plan rates are uniformly moderate (0.47--0.73) and
outcome-invariant.

% ===========================================================================
\paragraph{Action-conditioned entropy: tool use lowers output entropy but
does not guarantee correctness.}
% ===========================================================================

Table~\ref{tab:cross-action-entropy} reports mean per-token log-probability
entropy conditioned on the presence or absence of specific action patterns
in the trajectory.

\begin{table}[t]
\centering
\small
\caption{%
  Mean per-token entropy conditioned on action pattern presence, by harness
  and outcome on GPQA-Diamond.
  ``Has tool'' = trajectory contains at least one tool\_use event.
  ``Has recover'' = trajectory contains at least one recover event.
  ``Verify after answer'' = verify action occurs after the terminal answer.
  $\Delta$Ent = entropy(wrong) $-$ entropy(correct) for the same pattern.
}
\label{tab:cross-action-entropy}
\setlength{\tabcolsep}{3pt}
\begin{tabular}{llrrrr}
\toprule
Harness & Action pattern & Ent. (C) & Ent. (W) & $\Delta$Ent & $n$ (C\,/\,W) \\
\midrule
\multirow{4}{*}{OpenClaw}
  & Has tool use      & 0.216 & 0.264 & +0.048 & 29\,/\, 5 \\
  & No tool use       & 0.247 & 0.125 & $-$0.122 & 102\,/\,62 \\
  & Has recover       & 0.201 & 0.169 & $-$0.032 & 30\,/\, 7 \\
  & Verify after ans. & 0.242 & 0.148 & $-$0.094 & 74\,/\,39 \\
\midrule
\multirow{4}{*}{OpenCode}
  & Has tool use      & 0.251 & \textbf{0.378} & \textbf{+0.127} & 80\,/\, 8 \\
  & No tool use       & 0.307 & 0.273 & $-$0.034 & 65\,/\,45 \\
  & Has verify        & 0.293 & 0.280 & $-$0.013 & 86\,/\,26 \\
  & Has plan          & 0.289 & 0.265 & $-$0.024 & 44\,/\,16 \\
\midrule
\multirow{4}{*}{ZeroClaw}
  & Has verify        & 0.245 & 0.350 & \textbf{+0.105} & 90\,/\,21 \\
  & Has plan          & 0.239 & 0.352 & \textbf{+0.113} & 122\,/\,27 \\
  & Sustained reason  & 0.245 & 0.359 & \textbf{+0.114} & 147\,/\,34 \\
\bottomrule
\end{tabular}
\end{table}

\findingbox{Tool use is an entropy suppressor, not a correctness guarantee.
Under OpenCode, trajectories that invoke tools have lower entropy when
correct (0.25) but higher entropy when wrong (0.38), producing the only
positive $\Delta$Ent in the table.
The model persists in tool use under uncertainty, and that uncertainty is
visible in the output entropy.}

Table~\ref{tab:cross-action-entropy} reveals three regularities.

First, tool use lowers entropy on correct trajectories across both
tool-using harnesses.
OpenClaw correct trajectories with tool use have mean entropy 0.22 vs.\
0.25 without; OpenCode correct with tool use have 0.25 vs.\ 0.31 without.
Tool execution anchors the reasoning and the model's output becomes more
deterministic, which is appropriate when the tool result confirms the
reasoning path.

Second, tool use on wrong trajectories carries the opposite signal.
OpenCode wrong trajectories with tool use have entropy 0.38, which is
higher than any other cell in the table and 0.13 above correct
tool-using trajectories.
The model invokes tools under uncertainty, the tool output does not
resolve the uncertainty, and the model's subsequent reasoning retains
high entropy.
This is the failure mode the ARM analysis identifies as
``uncertainty without resolution'': the model uses tools exploratively
but cannot convert the output into a confident answer.

Third, ZeroClaw's action-conditioned entropy shows a uniform positive
$\Delta$Ent: wrong trajectories have higher entropy than correct ones
for every action pattern (has verify +0.11, has plan +0.11, sustained
reason +0.11).
This is the signature of genuine reasoning under uncertainty within a
bounded budget.
The model is uncertain, the entropy reflects that uncertainty, but the
template prevents the uncertainty from spiraling into repetition.
The entropy is high rather than collapsed, consistent with the context
rotting finding that ZeroClaw does not enter the catastrophic rotting
regime.
